\chapter{1977年北京卷（文）}

\section{解答题}

\begin{problem}
计算：\(3^{0}+3^{-1}-\left(1\frac{7}{9}\right)^{\frac{1}{2}}\).
\end{problem}

\begin{solution}
因为 \(3^{0}=1\)，\(3^{-1}=\frac{1}{3}\)，且 \(1\frac{7}{9}=\frac{16}{9}\)，所以 \(\left(1\frac{7}{9}\right)^{\frac{1}{2}}=\sqrt{\frac{16}{9}}=\frac{4}{3}\). 因此原式 \(=1+\frac{1}{3}-\frac{4}{3}=0\).
\end{solution}

\begin{problem}
化简：\(\frac{\sqrt{6}+\sqrt{2}}{\sqrt{6}-\sqrt{2}}\).
\end{problem}

\begin{solution}
分母有理化，得
\[
\frac{\sqrt{6}+\sqrt{2}}{\sqrt{6}-\sqrt{2}}=\frac{\left(\sqrt{6}+\sqrt{2}\right)^2}{6-2}=\frac{8+4\sqrt{3}}{4}=2+\sqrt{3}
\]
所以原式的值为 \(2+\sqrt{3}\).
\end{solution}

\begin{problem}
解方程：\(\frac{1}{x-1}+1=\frac{4x-2}{x^{2}-1}\).
\end{problem}

\begin{solution}
由分式有意义可知 \(x\neq1\)，且 \(x\neq-1\). 在此条件下，原方程两边同乘 \(x^{2}-1=(x-1)(x+1)\)，得 \(x+1+x^{2}-1=4x-2\)，即 \(x^{2}-3x+2=0\)，所以 \(\left(x-1\right)\left(x-2\right)=0\). 因此 \(x=1\) 或 \(x=2\). 由于 \(x=1\) 不满足原方程的定义域，舍去，故原方程的解为 \(x=2\).
\end{solution}

\begin{problem}
不查表求 \(\sin 105^\circ\) 的值.
\end{problem}

\begin{solution}
利用和角公式，得
\[
\sin105^\circ=\sin\left(60^\circ+45^\circ\right)=\sin60^\circ\cos45^\circ+\cos60^\circ\sin45^\circ=\frac{\sqrt{3}}{2}\cdot\frac{\sqrt{2}}{2}+\frac{1}{2}\cdot\frac{\sqrt{2}}{2}=\frac{\sqrt{6}+\sqrt{2}}{4}
\]
所以 \(\sin105^\circ=\frac{\sqrt{6}+\sqrt{2}}{4}\).
\end{solution}

\begin{problem}
一个正三棱柱形的零件，它的高是 \(10\mathrm{~cm}\)，底面边长是 \(2\mathrm{~cm}\)，求它的体积.
\end{problem}

\begin{solution}
正三棱柱的底面是边长为 \(2\mathrm{~cm}\) 的正三角形. 由等边三角形的高为 \(\sqrt{2^{2}-1^{2}}=\sqrt{3}\mathrm{~cm}\)，得底面积为 \(\frac{1}{2}\times2\times\sqrt{3}=\sqrt{3}\mathrm{~cm}^{2}\). 因此正三棱柱的体积为 \(V=\sqrt{3}\times10=10\sqrt{3}\mathrm{~cm}^{3}\).
所以它的体积为 \(10\sqrt{3}\mathrm{~cm}^{3}\).
\end{solution}

\begin{problem}
一条直线过点 \(\left(1,-3\right)\)，并且与直线 \(2x+y-5=0\) 平行，求这条直线的方程.
\end{problem}

\begin{solution}
所求直线与 \(2x+y-5=0\) 平行，因此可设其方程为 \(2x+y+c=0\). 将点 \(\left(1,-3\right)\) 代入，得 \(2-3+c=0\)，所以 \(c=1\). 故所求直线的方程为 \(2x+y+1=0\).
\end{solution}

\begin{problem}
证明：等腰三角形两腰上的高相等.
\end{problem}

\begin{solution}
设等腰三角形为 \(\triangle ABC\)，其中 \(AB=AC\). 从点 \(B\) 向腰 \(AC\) 作垂线，垂足为 \(D\)；从点 \(C\) 向腰 \(AB\) 作垂线，垂足为 \(E\). 则 \(BD\)、\(CE\) 分别是两腰上的高.

设三角形面积为 \(S\). 以 \(AC\) 和 \(AB\) 分别为底边时，三角形面积可表示为 \(S=\frac{1}{2}AC\cdot BD=\frac{1}{2}AB\cdot CE\).
因为 \(AB=AC\neq0\)，所以由上式可得 \(BD=CE\). 因而等腰三角形两腰上的高相等.
\end{solution}

\begin{problem}
为了测湖岸边 \(A\)，\(B\) 两点的距离，选择一点 \(C\)，测得 \(CA=50\text{ 米}\)，\(CB=30\text{ 米}\)，\(\angle ACB=120^\circ\)，求 \(AB\).
\end{problem}

\begin{solution}
在 \(\triangle ACB\) 中，由余弦定理，\(AB^{2}=CA^{2}+CB^{2}-2\,CA\cdot CB\cos\angle ACB\).
代入已知数据，得
\[
AB^{2}=50^{2}+30^{2}-2\times50\times30\cos120^\circ=2500+900+1500=4900
\]
因为长度为正，所以 \(AB=\sqrt{4900}=70\text{ 米}\).
\end{solution}

\begin{problem}
在 \(2\) 和 \(30\) 中间插入两个正数，这两个正数插入后使前三个数成等比数列，后三个数成等差数列，求插入的两个正数.
\end{problem}

\begin{solution}
设插入的两个正数依次为 \(x\)，\(y\)，则四个数依次为 \(2\)，\(x\)，\(y\)，\(30\). 由前三个数成等比数列，得 \(x^{2}=2y\). 由后三个数成等差数列，得 \(2y=x+30\). 联立两式，得到 \(x^{2}=x+30\)，即 \(x^{2}-x-30=0\)，所以 \(\left(x-6\right)\left(x+5\right)=0\).
因为 \(x\) 为正数，所以 \(x=6\)，进而 \(y=\frac{x+30}{2}=18\). 检验可知 \(2,6,18\) 成等比数列，\(6,18,30\) 成等差数列，故插入的两个正数为 \(6\)，\(18\).
\end{solution}

\begin{problem}
已知二次函数 \(y=x^{2}-4x+3\).

\begin{enumerate}
\item 求出它的图象的顶点坐标和对称轴方程；
\item 画出它的图象；
\item 求出它的图象与直线 \(y=x-3\) 的交点坐标.
\end{enumerate}
\end{problem}

\begin{solution}
\begin{enumerate}
\item 配方得 \(y=x^{2}-4x+3=\left(x-2\right)^{2}-1\)，
所以抛物线的顶点坐标为 \(\left(2,-1\right)\)，对称轴方程为 \(x=2\).

\item 由（1）知，图象是开口向上的抛物线，顶点为 \(\left(2,-1\right)\)，对称轴为 \(x=2\). 令 \(y=0\)，得 \(x^{2}-4x+3=0\)，即 \(\left(x-1\right)\left(x-3\right)=0\)，所以图象与 \(x\) 轴交于 \(\left(1,0\right)\) 和 \(\left(3,0\right)\). 又当 \(x=0\) 时 \(y=3\)，故图象经过 \(\left(0,3\right)\)；由关于直线 \(x=2\) 对称，还经过 \(\left(4,3\right)\). 据此描点并作出开口向上的抛物线，即得所求图象.

\item 联立抛物线与直线的方程，得
\[
\begin{cases}
y=x^{2}-4x+3,\\
y=x-3.
\end{cases}
\]
因此 \(x^{2}-4x+3=x-3\)，即 \(x^{2}-5x+6=0\)，所以 \(x=2\) 或 \(x=3\). 当 \(x=2\) 时，\(y=-1\)；当 \(x=3\) 时，\(y=0\). 故交点坐标为 \(\left(2,-1\right)\) 和 \(\left(3,0\right)\).
\end{enumerate}
\end{solution}
